\documentclass[a4paper,11pt]{article}
\usepackage{exam-style-842-unit}
\examinehead{842 信息化概论 · 第09单元练习册}{练习二 参考答案与解析}
\begin{document}

\unitanswerscover{842 信息化概论 · 单元练习册}{第09单元参考答案}{练习二（3S技术与农业专家系统）的答案与解析；每道判断题均给出改正后的正确表述，综合题另给出答题框架、给分点与分值分配、常见失分。}

\answerpart{一、单项选择题（10 题，每题 3 分，共 30 分）}

\q A。\begin{analysis}遥感依据电磁波理论：一切物体都在反射或辐射电磁波，不同地物（作物、土壤、水体、岩石）的光谱特性不同，传感器接收并记录这些差异，经处理解译即可识别地物及其状态。B 重量差异、C 酸碱度差异、D 价格差异都不是遥感探测的物理基础。\end{analysis}

\q A。\begin{analysis}归一化植被指数（NDVI）由近红外与红光波段的反射率组合而成，绿色植被在该指数上表现为高值，因而广泛用于作物长势监测、植被覆盖度估算和产量预测。B 土壤容重需田间实测；C 农药残留需化学分析；D 化肥价格指数属于市场信息，均不能由遥感直接获取。\end{analysis}

\q A。\begin{analysis}GIS 与一般数据库管理系统的根本区别在于它管理的是具有空间位置、拓扑关系和属性信息的地理数据，并能在空间维度上查询、叠加与分析。B、C 把 GIS 的功能窄化为存储文字或制图；D 恰恰相反，数据是 GIS 的核心组成之一。\end{analysis}

\q A。\begin{analysis}GIS 的核心功能是空间数据的采集与编辑、存储与管理、空间查询与分析、输出与可视化，其中空间分析（叠加、缓冲区、网络、地形分析等）最能体现其价值。B、C、D 分别是排版、视频、语音处理软件的功能。\end{analysis}

\q C。\begin{analysis}GPS 接收机测量到各卫星的伪距，未知量包括三维位置（经度、纬度、高度）和接收机钟差共 4 个，因此至少需要同时观测 4 颗卫星才能解算。只观测 3 颗卫星时无法同时消除钟差，A、B 不足；D 的 6 颗并非最低要求。\end{analysis}

\q A。\begin{analysis}GPS 由空间部分（卫星星座）、地面控制部分（主控站、监控站、注入站）和用户设备部分（接收机）三大部分组成（教材 6.1.2，印刷页 100）。B 是物联网的四层体系；C、D 是一般系统或信息处理流程的划分，不是 GPS 的组成。\end{analysis}

\q A。\begin{analysis}北斗卫星导航系统在提供定位、导航、授时服务的同时，具有短报文通信能力，这是其区别于 GPS 的特色功能（教材 6.1.2，印刷页 100）。B 北斗同样可导航；C 北斗可授时且精度较高；D 北斗是军民两用的开放系统。\end{analysis}

\q B。\begin{analysis}3S 集成中，遥感负责获取面状信息，GPS／北斗负责提供精确位置，GIS 则是空间数据的“管理平台”和“分析大脑”，承担数据管理、叠加分析、处方图生成等任务。A 遥感是信息源；C 定位系统提供位置；D 虚拟现实用于模拟展示。\end{analysis}

\q A。\begin{analysis}农业专家系统由知识库、推理机、综合数据库、人机接口和知识获取等部分组成，其中知识库存放专家知识、推理机负责推理，二者是系统的核心。B 是输入输出设备；C 是网络硬件；D 是通用的数据库软件，都不构成专家系统的核心。\end{analysis}

\q B。\begin{analysis}农业数据库及管理系统用于农业资源的组织、存储、检索、维护与共享，可管理品种资源、土壤、气象、病虫害、市场等多类数据。A 农业数据还包括数值、图形图像、时序等多类型；C 数据库正是为农业生产与管理决策服务；D 没有数据库管理系统就无法有效组织与维护数据。\end{analysis}

\answerpart{二、多项选择题（4 题，每题 4 分，共 16 分）}

\q ABC。\begin{analysis}遥感可大范围、快速、周期性地获取作物长势、种植面积、产量与灾害信息，是农业监测的重要信息源。D 错误：施肥作业是变量作业装备的任务，遥感只能提供“哪里需要施肥”的信息，不能直接完成施肥。\end{analysis}

\q ABCD。\begin{analysis}GIS 的基本功能包括数据采集与编辑、数据存储与管理、空间查询与空间分析、数据输出与可视化，并可通过二次开发嵌入专业模型。四项全对。\end{analysis}

\q ABC。\begin{analysis}农业专家系统一般由知识库、推理机、综合数据库、人机接口和知识获取五部分组成，有的把解释机构单列。D 错误：操作系统是计算机的系统软件，不是专家系统的功能组成。\end{analysis}

\q ABCD。\begin{analysis}北斗的高精度定位与导航能力可支撑农机自动驾驶、精准播种、变量施肥、植保无人机导航、联合收割机测产和农田测绘，其短报文通信还可用于边远地区作业监控与调度。四项全对。\end{analysis}

\answerpart{三、判断题（5 题，每题 2 分，共 10 分）}

\q ×。改正：遥感分为被动遥感和主动遥感，被动遥感依赖地物反射太阳光，主动遥感（如合成孔径雷达）自身发射电磁波并接收回波，可全天时、全天候工作，因此遥感并非夜间和阴天完全不能工作。\begin{analysis}把遥感等同于“光学被动遥感”，忽视了微波遥感等主动遥感手段，是常见误解。\end{analysis}

\q √。\begin{analysis}GIS 的叠加分析可把土壤、地形、土地利用、气象等图层叠合，按空间位置生成新的属性组合，是土地适宜性评价、选址与区划的常用方法。\end{analysis}

\q ×。改正：由于接收机钟差未知，GPS 定位至少需要同时观测 4 颗卫星，才能解算经度、纬度和高度并校正钟差；3 颗卫星只能解算三维位置而无法处理钟差。\begin{analysis}“3 颗卫星定位”是常见错误记忆，考试中常作为判断题或单选题的考点。\end{analysis}

\q ×。改正：农业专家系统依赖知识库中的专家知识进行推理，必须通过知识获取模块不断补充、更新知识，且难以完全替代人类专家处理新问题。\begin{analysis}专家系统只能在知识覆盖的范围内给出建议，知识获取是其主要“瓶颈”，因此“完全替代人类专家、无需更新知识”的说法错误。\end{analysis}

\q √。\begin{analysis}虚拟农业以虚拟现实、仿真和作物模型构建虚拟农场，可对品种、播期、肥水方案进行模拟评价，并用于农技培训与园区规划，降低田间试验成本。\end{analysis}

\answerpart{四、配对题（6 对，每对 1 分，共 6 分）}

\q （1）—B\quad （2）—A\quad （3）—C\quad （4）—D\quad （5）—E\quad （6）—F。
\begin{analysis}遥感负责获取地物光谱信息；GIS 负责空间数据管理与空间分析；北斗提供三维定位、导航并具有短报文通信功能；知识库存放领域专家知识与规则；推理机依据事实与规则推出结论；虚拟植物实现植物形态与生长过程的可视化模拟。（对应关系见本单元知识地图与术语辨析。）\end{analysis}

\answerpart{五、简答题（4 题，每题 6 分，共 24 分）}

\q 参考答案要点：
（1）物理基础（2 分）：一切物体都在反射或辐射电磁波，不同地物具有不同的光谱特性；传感器接收、记录地物电磁波信息，经处理与解译获取地物性质、状态与变化。
（2）技术环节（2 分）：遥感平台（卫星、飞机、无人机）搭载传感器获取数据 $\rightarrow$ 数据传输与预处理（辐射校正、几何校正）$\rightarrow$ 图像处理与解译（分类、指数计算）$\rightarrow$ 应用分析。
（3）特点（1 分）：探测范围大、获取信息快、周期短、受地面条件限制少，可反映动态变化。
（4）农业应用（1 分，答出 2 点即可）：作物长势监测与产量估测（NDVI 等植被指数）、种植面积与种植结构调查、土壤墒情与旱情监测、病虫害与灾害监测、农业资源调查。

\q 参考答案要点：
（1）组成（2 分）：硬件（计算机、输入输出设备）、软件（GIS 平台软件与应用软件）、数据（空间数据与属性数据，是核心）、人员（操作与管理人员）、方法或规范（空间分析模型与标准）。
（2）基本功能（2 分）：数据采集与编辑、数据存储与管理、空间查询与空间分析（叠加分析、缓冲区分析、网络分析、地形分析）、数据输出与可视化、二次开发。
（3）农业应用（2 分，答出 2 点即可）：土壤与土地利用管理；土地适宜性评价与农业区划；精准农业中划分管理单元、生成处方图与产量分布图；农业灾害评估与灾后损失统计；农田水利与道路规划。

\q 参考答案要点：
（1）组成（2 分）：GPS 由空间部分（卫星星座）、地面控制部分（主控站、监控站、注入站）、用户设备部分（接收机）组成；北斗由空间段、地面段、用户段组成（教材 6.1.2，印刷页 100—101）。
（2）定位原理（2 分）：卫星发播含星历与时间信息的导航信号，接收机测量到各卫星的伪距，根据“距离交会”原理，至少同时观测 4 颗卫星，即可解算接收机的经度、纬度、高度和钟差；北斗在此基础上还可提供短报文通信。
（3）农业应用（2 分，答出 2 点即可）：农机自动驾驶与辅助导航；精准播种、变量施肥、变量喷药；联合收割机测产与作业监控；农田测绘与土地平整；植保无人机导航与畜牧、渔船定位管理。

\q 参考答案要点：
（1）组成（3 分）：知识库（领域专家的知识、经验与规则）、推理机（推理策略与控制程序）、综合数据库（初始事实、用户输入与中间结果）、人机接口（人机交互界面）、知识获取（知识录入、维护与更新）；解释机构可单列，用于说明推理依据。教材 10.5.2，印刷页 195—197，从人工智能角度把专家系统概括为知识库、数据库、推理机构、解释机制和知识获取五部分，与本处口径一致，只是“解释机制”单列、“综合数据库”对应“数据库”。
（2）工作原理（3 分）：用户通过人机接口输入地块、作物、环境等事实 $\rightarrow$ 事实存入综合数据库 $\rightarrow$ 推理机按正向或反向推理匹配知识库中的规则 $\rightarrow$ 推出结论并给出施肥、灌溉、防治等建议 $\rightarrow$ 解释机构说明推理依据，知识获取模块根据反馈补充更新知识库。
（3）典型应用（可并入第（3）点）：病虫害诊断与防治、测土配方施肥、灌溉与栽培管理决策、品种选择、畜禽疾病诊断等。

\answerpart{六、综合题（1 题，共 14 分）}

\q 参考答案与答题框架：

\textbf{【破题句】}该体系以 RS 为“天”、无人机与地面物联网为“空—地”信息源，以 GIS 为空间分析与决策平台，以北斗为作业定位与导航手段，形成“监测—分析—作业”闭环的 3S 集成应用，并把农业专家系统嵌入决策环节，实现从“看得见”到“算得准”再到“作业精”的转变。

\textbf{（1）3S 的任务与集成链条（6 分）}
\begin{analysis}
RS 任务（2 分）：卫星遥感和无人机巡查获取作物分布、长势、病虫害与灾情等面状、动态信息，回答“哪里、发生了什么”；GIS 任务（2 分）：管理土壤、气象、产量等空间图层，叠加分析划分管理单元并生成处方图，回答“为什么、怎么办”；北斗任务（1 分）：提供厘米级至米级位置与导航，引导农机按处方图精确作业，回答“精确到哪里作业”。集成链（1 分）：RS 获取信息 $\rightarrow$ GIS 处理分析、生成处方 $\rightarrow$ 北斗定位导航、驱动变量作业 $\rightarrow$ 作业结果反馈回 GIS 与遥感，形成闭环。
\end{analysis}

\textbf{（2）农业专家系统的工作过程（4 分）}
\begin{analysis}
按四步给分，每步 1 分：
（1）事实输入：通过人机接口输入地块位置、土壤养分、作物品种与长势、气象与病虫信息，存入综合数据库；
（2）知识匹配：推理机读取综合数据库中的事实，检索知识库中的规则（如“pH 偏低且缺磷 $\rightarrow$ 施石灰并增施磷肥”）；
（3）推理结论：按正向推理得出结论，给出施肥种类用量、防治时期与药剂等建议；
（4）解释与更新：解释机构说明推理依据，农技人员反馈实际效果，知识获取模块补充、修正知识库。
\end{analysis}

\textbf{（3）作用与价值（4 分）}
\begin{analysis}
答出 4 点、每点 1 分，言之成理即可：
（1）节本增效：按需施肥施药灌溉，减少化肥农药与水肥浪费，降低生产成本；
（2）提高决策科学性：把专家经验固化为可复制的知识，缓解基层农技人员不足；
（3）提升监测预警能力：及早发现病虫害、旱情与灾情，减少损失；
（4）支撑绿色与质量安全：减少面源污染，结合作业记录支撑农产品质量追溯；
（5）促进规模经营与服务化：为合作社、社会化服务组织提供标准化作业能力，推动小农户与现代农业衔接。
\end{analysis}

\textbf{【给分点与分值分配】}（1）6 分：RS 2 分、GIS 2 分、北斗 1 分、集成链 1 分；（2）4 分：事实输入、知识匹配、推理结论、解释更新各 1 分；（3）4 分：作用价值 4 点。

\textbf{【常见失分】}一是把 3S 的分工写反，如让 GPS 做空间分析、让 GIS 获取遥感影像，技术定位不清；二是只写“3S 集成就是三个技术结合”，没有写出“遥感发现—GIS 分析—北斗定位作业”的链条与反馈闭环；三是专家系统只答组成不答工作过程，或把人机接口、知识获取漏掉；四是作用价值空泛，不结合节本增效、绿色生产、农技推广等具体场景；五是全篇不结合“天空地”监测与水稻、柑橘等具体作物，论述缺乏落点。

\end{document}
